\documentclass[11pt,a4paper]{article}

\usepackage[utf8]{inputenc}
\usepackage[T1]{fontenc}
\usepackage{lmodern}
\usepackage[brazilian]{babel}
\usepackage[margin=2.5cm]{geometry}
\usepackage{titlesec}
\usepackage{enumitem}
\usepackage{xcolor}
\usepackage{fancyhdr}
\usepackage{hyperref}
\usepackage{booktabs}
\usepackage{array}
\usepackage{parskip}
\usepackage{microtype}
\sloppy

\definecolor{accent}{HTML}{1F6E66}
\definecolor{accentwarm}{HTML}{A8672A}
\definecolor{inksoft}{HTML}{444444}

\titleformat{\section}
  {\normalfont\Large\bfseries\color{accent}}
  {\thesection}{0.6em}{}
\titleformat{\subsection}
  {\normalfont\large\bfseries\color{black}}
  {\thesubsection}{0.6em}{}

\hypersetup{
  colorlinks=true,
  linkcolor=accent,
  urlcolor=accent,
  citecolor=accent,
  pdftitle={Roteiro Teórico - Rede Social Simplificada},
  pdfauthor={Estruturas de Dados - UFU}
}

\pagestyle{fancy}
\fancyhf{}
\renewcommand{\headrulewidth}{0.4pt}
\fancyhead[L]{\small\textcolor{inksoft}{Rede Social Simplificada --- Roteiro Teórico}}
\fancyhead[R]{\small\textcolor{inksoft}{\thepage}}

\setlist[itemize]{leftmargin=1.4em, itemsep=2pt, topsep=2pt}
\setlist[enumerate]{leftmargin=1.6em, itemsep=2pt, topsep=2pt}

\newcommand{\slideref}[1]{\textcolor{accentwarm}{\texttt{[Slide #1]}}}
\newcommand{\code}[1]{\texttt{#1}}

\title{
  \vspace{-2cm}
  {\Large\textcolor{inksoft}{Estruturas de Dados --- UFU --- Prof. Camila Davi Ramos --- 2026/1}}\\[0.6cm]
  {\huge\bfseries\color{accent} Roteiro Teórico da Apresentação}\\[0.3cm]
  {\Large Rede Social Simplificada --- Projeto Final, Tema 04}
}
\author{}
\date{}

\begin{document}
\maketitle
\thispagestyle{fancy}

\noindent\textbf{Como usar este documento:} ele acompanha os 14 slides de
\texttt{slides.html} e o roteiro de demonstração ao vivo em
\texttt{roteiro\_demo.md}. Enquanto aquele roteiro foca no \emph{o que digitar
e o que vai aparecer na tela}, este aqui é o material de apoio teórico --- os
conceitos de estruturas de dados e algoritmos por trás de cada slide,
resumidos para serem falados em 30--60 segundos cada. É o conteúdo que
justifica, na introdução do relatório, por que cada estrutura foi escolhida.

\tableofcontents
\vspace{0.4cm}

%=================================================================
\section{Introdução: o que é um TAD e por que isso importa}
\slideref{1-2}

Um \textbf{Tipo Abstrato de Dados (TAD)} é a definição de um conjunto de
dados junto com as operações permitidas sobre eles, \emph{sem} amarrar essa
definição a uma implementação específica. Uma pilha, por exemplo, é definida
pelo comportamento ``último a entrar, primeiro a sair'' e pelas operações
\code{empilhar}/\code{desempilhar} --- não importa se por baixo dos panos ela
usa um vetor ou uma lista encadeada. Essa separação entre \emph{interface} e
\emph{implementação} é o que permite trocar a forma como uma estrutura
funciona por dentro sem quebrar quem a utiliza por fora --- e é exatamente
por isso que o projeto exige cada estrutura como um par independente
\code{arquivo.h} (a interface, o contrato) e \code{arquivo.c} (a
implementação).

A escolha de qual estrutura usar em cada parte do sistema não é arbitrária:
depende de \textbf{qual operação precisa ser rápida}. Medimos isso com
\textbf{notação O grande (Big-O)}, que descreve como o tempo (ou espaço)
necessário cresce conforme o tamanho da entrada $n$ aumenta:

\begin{itemize}
  \item $O(1)$ --- tempo constante, independe de $n$ (ex.: acessar um bucket de hash).
  \item $O(\log n)$ --- cresce devagar; dobrar $n$ só soma um passo a mais (ex.: buscar em uma árvore balanceada).
  \item $O(n)$ --- cresce proporcional a $n$ (ex.: percorrer uma lista inteira).
  \item $O(V+E)$ --- para grafos, cresce com o número de vértices $V$ mais o de arestas $E$ (ex.: BFS).
\end{itemize}

O projeto usa seis TADs justamente porque cada funcionalidade obrigatória tem
um perfil de acesso diferente: buscar por id pede $O(\log n)$ garantido,
validar login duplicado pede $O(1)$ médio, listar publicações pede inserção
$O(1)$, e assim por diante --- é o que as próximas seções detalham.

%=================================================================
\section{Arquitetura em fachada}
\slideref{3}

O sistema segue o \textbf{padrão de projeto Fachada (Facade)}: um único
módulo (\code{rede\_social.c}) concentra todo o estado (a árvore, as duas
tabelas hash, o grafo, a pilha de histórico) e expõe uma função de alto nível
por funcionalidade obrigatória (\code{rs\_cadastrar\_usuario},
\code{rs\_adicionar\_amizade}, etc.). O \code{main.c} não sabe --- e não
precisa saber --- como cada operação é resolvida por dentro; ele só faz o
loop do menu e chama a fachada.

Essa separação entre \emph{lógica de domínio} e \emph{interface de usuário} é
o que permite testar todas as regras de negócio (duplicidade de login,
amizade repetida, ordenação de sugestões) sem precisar simular um terminal
--- é exatamente o que \code{make test} faz.

%=================================================================
\section{As seis estruturas de dados obrigatórias}
\slideref{4}

Esta é a seção central: para cada estrutura, o conceito, como ela resolve
colisões/desbalanceamento (quando aplicável) e por que ela foi a escolha
certa para aquele papel no sistema.

%-----------------------------------------------------------------
\subsection{Árvore AVL --- índice por id}

Uma \textbf{árvore binária de busca (BST)} organiza chaves de forma que, em
cada nó, tudo à esquerda é menor e tudo à direita é maior. Isso permite
busca, inserção e remoção em tempo proporcional à \textbf{altura} da árvore.
O problema: se os valores chegarem em ordem crescente (ou decrescente), uma
BST comum degenera em uma lista encadeada disfarçada --- altura $O(n)$, e a
promessa de busca rápida desaparece exatamente no pior momento.

A \textbf{árvore AVL} (Adelson-Velsky e Landis, 1962) resolve isso
acrescentando uma invariante: para todo nó, a diferença de altura entre as
subárvores esquerda e direita (o \emph{fator de balanceamento}) nunca pode
ser maior que 1. Sempre que uma inserção quebra essa invariante em algum
ancestral, uma \textbf{rotação} (simples ou dupla) reorganiza um pequeno
trecho da árvore para restaurar o balanceamento em tempo $O(1)$. O resultado
é uma altura sempre $O(\log n)$, garantida --- não apenas no caso médio.

\textbf{Por que aqui:} índice \code{id -> Usuario*}. Ids são gerados de
forma sequencial (1, 2, 3, ...) --- exatamente o padrão de entrada que
quebraria uma BST comum. A AVL garante busca por id em $O(\log n)$ mesmo
nesse cenário adversarial.

%-----------------------------------------------------------------
\subsection{Tabela hash --- índice por login e por nome}

Uma \textbf{tabela hash} usa uma \textbf{função hash} $h(chave)$ para
calcular diretamente a posição (bucket) de um vetor onde aquele dado deve
estar, trocando a busca por comparações sucessivas (como numa árvore) por um
cálculo direto --- em média, $O(1)$. Uma boa função hash distribui as chaves
uniformemente pelos buckets; o projeto usa a \textbf{djb2}, um hash
multiplicativo simples e rápido para strings.

Como duas chaves diferentes podem calcular o mesmo bucket
(\textbf{colisão}), é preciso uma estratégia de resolução. O projeto usa
\textbf{encadeamento separado}: cada bucket é, ele mesmo, uma lista
encadeada de entradas --- colisões viram apenas uma lista um pouco mais
longa para percorrer, sem invalidar outras posições do vetor.

O desempenho $O(1)$ depende do \textbf{fator de carga} (número de elementos
dividido pelo número de buckets) se manter baixo. Quando ultrapassa um
limiar (aqui, 0,75), a tabela faz um \textbf{rehash}: cria um vetor maior e
redistribui todas as chaves --- mantendo o custo médio amortizado em
$O(1)$ mesmo com a tabela crescendo indefinidamente.

\textbf{Por que aqui:} duas instâncias independentes. Uma indexa por
\textbf{login} (chave única, usada para barrar cadastro duplicado em
$O(1)$); outra indexa por \textbf{nome} (chave repetível, cada bucket
acumula uma lista de usuários com aquele nome, usada na busca por nome).

%-----------------------------------------------------------------
\subsection{Lista encadeada --- publicações e adjacência}

Uma \textbf{lista encadeada} é uma sequência de nós onde cada nó guarda um
dado e um ponteiro para o próximo --- ao contrário de um vetor, os elementos
não precisam ocupar posições contíguas de memória. Isso torna a
\textbf{inserção no início} uma operação $O(1)$: basta criar o novo nó
apontando para o antigo início, sem deslocar nada. O custo é o acesso
aleatório, que passa a ser $O(n)$ (é preciso percorrer nó a nó).

\textbf{Por que aqui:} cada usuário guarda suas publicações numa lista
encadeada própria. Inserir sempre no início faz a publicação mais recente
aparecer primeiro \emph{de graça}, sem precisar ordenar nada na hora de
exibir. A mesma estrutura genérica (dado + destrutor) é reaproveitada nas
listas de adjacência do grafo, guardando os vizinhos de cada vértice.

%-----------------------------------------------------------------
\subsection{Fila --- motor da busca em largura}

Uma \textbf{fila} segue a disciplina \textbf{FIFO} (\emph{First In, First
Out}): o primeiro elemento a entrar (\code{enfileirar}) é o primeiro a sair
(\code{desenfileirar}). É a estrutura de controle por trás da
\textbf{busca em largura (BFS)}, detalhada na Seção~\ref{sec:bfs}.

\textbf{Por que aqui:} sem uma fila real (FIFO), não há como garantir que o
grafo seja visitado \emph{nível por nível} durante a busca de conexão e de
sugestões de amizade --- usar uma pilha ali, por exemplo, transformaria a
busca em profundidade (DFS), que não garante encontrar primeiro os vizinhos
mais próximos.

%-----------------------------------------------------------------
\subsection{Pilha --- histórico de atividades}

Uma \textbf{pilha} segue a disciplina oposta, \textbf{LIFO} (\emph{Last In,
First Out}): o último elemento inserido (\code{empilhar}) é o primeiro a
sair (\code{desempilhar}). Perguntar ``o que aconteceu por último?'' é uma
operação natural em $O(1)$ numa pilha --- basta olhar o topo.

\textbf{Por que aqui:} o histórico exige exibir as atividades da mais
recente para a mais antiga --- exatamente a ordem que uma pilha devolve ao
ser percorrida do topo para a base. Com capacidade fixa ($N=10$), a própria
pilha descarta o elemento mais antigo (a base) a cada novo evento acima da
capacidade, sem precisar de nenhuma lista auxiliar de controle.

%-----------------------------------------------------------------
\subsection{Grafo --- rede de amizades}

Um \textbf{grafo} é um conjunto de \textbf{vértices} (aqui, usuários) ligados
por \textbf{arestas} (aqui, amizades). Amizade é uma relação simétrica --- se
A é amigo de B, B é amigo de A --- então o grafo é \textbf{não dirigido}.

Duas formas clássicas de representar um grafo:
\begin{itemize}
  \item \textbf{Matriz de adjacência:} vetor $V \times V$ indicando se existe
    aresta entre cada par. Checar uma aresta é $O(1)$, mas o espaço é sempre
    $O(V^2)$, mesmo que a rede tenha poucas conexões.
  \item \textbf{Lista de adjacência:} cada vértice guarda apenas a lista dos
    seus vizinhos reais. Espaço $O(V+E)$ --- muito mais eficiente quando o
    grafo é \emph{esparso}, como é o caso de redes sociais reais (a maioria
    das pessoas não é amiga da maioria das outras).
\end{itemize}

O projeto usa lista de adjacência (reaproveitando o TAD Lista Encadeada por
vértice), pelo mesmo motivo que qualquer rede social real usa: o número de
amizades por pessoa é pequeno comparado ao total de usuários.

%=================================================================
\section{Busca em largura (BFS): o algoritmo por trás de duas funcionalidades}
\label{sec:bfs}
\slideref{9-10}

A \textbf{busca em largura} (\emph{Breadth-First Search}) percorre um grafo
a partir de um vértice de origem, visitando primeiro todos os vizinhos à
\textbf{distância 1}, depois todos à distância 2, e assim por diante ---
nunca avançando para um nível mais distante antes de esgotar o nível atual.
O algoritmo usa duas estruturas auxiliares:

\begin{enumerate}
  \item Uma \textbf{fila} (FIFO) com os vértices a visitar --- é o que impõe
    a ordem ``nível por nível'': o vértice enfileirado há mais tempo é
    sempre o próximo a ser expandido.
  \item Um vetor (ou conjunto) de \textbf{visitados}, para nunca processar o
    mesmo vértice duas vezes --- essencial em um grafo com ciclos.
\end{enumerate}

Complexidade $O(V+E)$: cada vértice entra na fila no máximo uma vez, e cada
aresta é examinada no máximo uma vez (uma de cada lado).

\textbf{Funcionalidade IX --- verificar conexão:} BFS a partir do usuário A,
parando assim que o usuário B for encontrado (conectados) ou a fila
esvaziar sem encontrá-lo (não conectados).

\textbf{Funcionalidade X --- sugestão de amizades:} a mesma BFS, mas
limitada a \textbf{distância 2} --- ou seja, ``amigos de amigos'' que ainda
não são amigos diretos. Para cada candidato encontrado, conta-se quantos
amigos em comum ele tem com o usuário consultado (interseção das listas de
adjacência dos dois), e a lista final é ordenada por
\begin{enumerate}
  \item número de amigos em comum, decrescente;
  \item em caso de empate, nome em ordem alfabética;
\end{enumerate}
cortada em no máximo 3 sugestões, como pede o enunciado.

%=================================================================
\section{Requisitos não funcionais: persistência e tratamento de erros}
\slideref{11}

\textbf{Persistência} é o princípio de que o estado de um programa não deve
desaparecer quando o processo termina. Na prática, isso significa
\textbf{serializar} as estruturas em memória (a árvore, as tabelas hash, o
grafo, as listas, a pilha) para um formato plano em arquivo ao encerrar, e
\textbf{desserializar} --- reconstruir cada estrutura repovoando-a a partir
do arquivo --- ao iniciar. O projeto usa um formato texto simples, com seções
delimitadas para usuários, amizades, publicações e histórico.

\textbf{Tratamento defensivo de erros} é o princípio de nunca assumir que
uma entrada externa (digitada pelo usuário, ou lida de um arquivo) é
válida. Toda operação da fachada retorna um código de sucesso/falha e
preenche uma mensagem de erro específica para o chamador decidir o que
fazer --- em vez de travar o programa ou continuar com um estado
inconsistente. É o princípio de ``falhar graciosamente'' (\emph{fail
gracefully}): o programa informa o problema e continua rodando.

%=================================================================
\section{Testes: unitário vs. integração}
\slideref{12}

Um \textbf{teste unitário} verifica uma única unidade de código
isoladamente --- por exemplo, testar só a árvore AVL, sem depender de
tabela hash, grafo ou qualquer outra parte do sistema. É o que os binários
\code{test\_avl}, \code{test\_hash\_tabela}, \code{test\_lista},
\code{test\_fila}, \code{test\_pilha} e \code{test\_grafo} fazem, cada um
focado em um TAD.

Um \textbf{teste de integração} verifica se várias unidades, combinadas,
produzem o comportamento esperado quando trabalham juntas --- é o papel de
\code{test\_rede\_social}, que exercita a fachada completa (cadastro, hash
de login, grafo de amizades e pilha de histórico atuando em conjunto) do
jeito que o sistema realmente é usado.

A combinação das duas camadas dá confiança em dois níveis: cada estrutura
funciona corretamente por si só, \textbf{e} a composição delas também
funciona.

%=================================================================
\section{Tabela-resumo: estrutura, operação-chave e complexidade}
\slideref{4}

\begin{center}
\small
\renewcommand{\arraystretch}{1.2}
\begin{tabular}{@{}
  >{\raggedright\arraybackslash}p{2.6cm}
  >{\raggedright\arraybackslash}p{4.4cm}
  >{\raggedright\arraybackslash}p{2.5cm}
  >{\raggedright\arraybackslash}p{4.2cm}@{}}
\toprule
\textbf{Estrutura} & \textbf{Operação-chave} & \textbf{Complexidade} & \textbf{Papel} \\
\midrule
Árvore AVL & Buscar/inserir por id & $O(\log n)$ & Índice \code{id -> Usuario} \\
Tabela hash & Buscar/inserir por login ou nome & $O(1)$ médio & Login único; busca por nome \\
Lista encadeada & Inserir no início & $O(1)$ & Publicações; adjacência do grafo \\
Fila & Enfileirar/desenfileirar & $O(1)$ & Motor da BFS \\
Pilha & Empilhar/desempilhar & $O(1)$ & Histórico (N=10 mais recentes) \\
Grafo (lista adj.) & BFS completa & $O(V+E)$ & Conexão e sugestão de amizades \\
\bottomrule
\end{tabular}
\end{center}

%=================================================================
\section*{Fechamento}
\addcontentsline{toc}{section}{Fechamento}
\slideref{13-14}

Todo esse arcabouço teórico existe para resolver um problema concreto: uma
rede social com cadastro, amizades, publicações e sugestões. A partir daqui,
o roteiro de demonstração ao vivo (\texttt{roteiro\_demo.md}) mostra cada
uma dessas estruturas em ação, com uma rede já populada de 10 usuários.

\end{document}
